% Folder 74, batch 1 (archivists' pages 1-20).
% Pass: Opus 5.5 (claude-opus-5-5), 2026-10-03 — first pass, unchecked against the pages by a human.
%
% Pages transcribed: 1, 3, 5, 6, 8, 10, 12, 13, 14, 15, 16, 18, 20.
% Pages absent: 2, 4, 7, 9, 11, 17, 19 (blank, or no mathematics).
%
% Fast drafting hand on loose sheets: the formulas carry, much of the
% connective prose does not. Notation fixed for the batch: I a set of
% cardinality 3, I = {i, j, k}; (E_i) a family of sets, Γ ⊂ ∏ E_i a ternary
% relation; (G_i) a family of groups acting on the E_i as bitorsors;
% δ the diagonal action; quasigroup-style identities in x, y, z on pp. 15-16.
\documentclass[11pt,a4paper]{article}
\input{../preamble/grothendieck}

\folder{74}
\batch{1}
\pages{1}{20}
\foldertitle{Complexes cubiques : notes manuscrites (s.d.).}
\dating{[à partir de 1976]}
\watermark{Édition de démonstration}

\begin{document}

\section{Relations ternaires et bitorseurs}

\page{1}
\textbf{I.} $I$ ens.\ de card~$3$.

$(E_i)_{i\in I}$ famille d'ensembles

$\Gamma\subset\prod_{i\in I}E_i$

Axiome : si \struck{\ill{}} $i,j,k\in I$ distincts (donc $I=\{i,j,k\}$) alors $\forall x_i\in E_i$, $x_j\in E_j$, $\exists!\,x_k\in E_k$ tel que $(x_i,x_j,x_k)\in\Gamma$.

Catégorie $\mathcal{C}$ ($I$ variable \uncertain{;} pour \uncertain{iso sur} $I$)

\textbf{II.} $I$ ens.\ de card~$3$
\drawing{Dans la marge gauche, un triangle aux sommets $E_i$, $E_j$, $E_k$, dont les côtés portent $G_k$, $G_j$, $G_i$.}

$(G_i)_{i\in I}$ famille de groupes

\struck{$E_{ij}$ : $L(G_i)\to L(G_j)$ syst.\ \uncertain{transitif} d'isom.\ \uncertain{extérieurs}} entre ces groupes
\note{la fin de ligne « entre ces groupes », écrite sous la ligne biffée, n'est pas biffée.}

$(E_i)_{i\in I}$ famille d'ens.

$\forall i\in I$, les groupes $\prod_{j\in I\setminus\{i\}}G_j$ opèrent sur $E_i$, de façon \add{(si $I=\{i,j,k\}$)} que $E_i$ deviennent \uncertain{des} \emph{bitorseurs} sous $G_j$, $G_k$. On peut \uncertain{renverser} l'opération de $G_\ell$ en l'opération \uncertain{inverse}.

Alors
\[
  \prod_{i\in I}\ \prod_{j\in I\setminus\{i\}}G_j \;=\; \prod_{j\in I}G_j^{\,I\setminus\{j\}}
\]
opère sur $P=\prod E_i$, donc $\prod_{j\in I}G_j$ opère \emph{diagonalement} \add{par $\delta$}. Soit $x=(x_i)_{i\in I}\in P$, alors
\[
  \begin{array}{cccc}
    E_i\times E_j\times E_k & x_i, & x_j, & x_k \;=\;\xi\\
    G_j\ G_k\ G_i\ G_j & g_jx_ig_k^{-1}, & g_kx_jg_i^{-1}, & g_ix_kg_j^{-1}
  \end{array}
\]
\[
  (g_jx_i,\ x_j,\ x_kg_j) \;=\; (x_ig_k^{-1},\ g_kx_jg_i^{-1},\ g_ix_k)
\]
\note{le bloc ci-dessus est entouré d'un trait qui le sépare du reste ; la seconde ligne de droite est en partie surchargée et sa lecture est douteuse.}
si $I=\{i,j,k\}$, les conditions \struck{\ill{}} \add{équivalentes} :
\note{la ligne d'introduction est surchargée de deux mots biffés et d'ajouts interlinéaires ; seul « équivalentes » se lit sûrement.}
\begin{itemize}
  \item[a$_1$)] $\delta(G_i)x\subset\delta(G_j\times G_k).x$
  \item[a$_2$)] $\delta(G_j)x\subset\delta(G_k\times G_i).x$
  \item[a$_3$)] $\delta(G_k)x\subset\delta(G_i\times G_j).x$
\end{itemize}

\struck{\ill{}} On peut encore les exprimer de la façon suivante : $\forall$ couple $(j,k)$ \struck{d'}$j,k\in I$, $j\neq k$, on \ill{}
\[
  \tau_{k,j}: E_i\to\mathrm{Isom}(G_j,G_k)
\]
[l'image \ill{} \ill{}
\note{la phrase se poursuit hors de la page ; la page 2 ne porte pas de mathématiques.}

\page{3}
dans \struck{l'\ill{}} \add{\uncertain{extérieurs}} \uncertain{ex. \ill{}}] et la condition \ill{}
\note{début de page : suite de la phrase ouverte au bas de la page 1 (« l'image … »).}
\drawing{À gauche, un premier triangle biffé, aux côtés marqués $G_i$, $G_j$, $G_k$.}

\begin{tikzcd}[column sep=large, row sep=large]
G_k \arrow[dr, "\tau_{i,k}(x_j)"'] & & G_j \arrow[ll, "\tau_{k,j}(x_i)"'] \\
& G_i \arrow[ur, "\tau_{j,i}(x_k)"'] &
\end{tikzcd}

\[
  \tau_{j,i}(x_k)\,\tau_{i,k}(x_j)\,\tau_{k,j}(x_i)=\mathrm{id}_{G_j}
\]
(condition qui \ill{} bien \uncertain{invariante} par \uncertain{permutations circulaires})

L'\emph{existence} d'un tel $x=(x_i)_{i\in I}$ équivaut à ce que les isom.\ \emph{extérieurs}
\[
  \tilde\tau_{ij}=\tilde\tau_{ij}(x_k)\quad \text{\struck{$(i,j,k\in I)$}}\quad i,j,k\in I \text{ distincts}
\]
définis par les $E_k$ ($k\in I$) entre les $G_i$ forment un \uncertain{système} \uncertain{transitif} d'isom.\ \uncertain{extérieurs}.

La condition $\mathrm{a}_i$ \uncertain{signifie} que si on \ill{} $x$ on \uncertain{dépend} \ill{} \ill{} \ill{} $\prod_{i\in I}G_i$, \struck{\ill{}} \add{\ill{}} \ill{} ainsi \ill{} \ill{} \uncertain{quelconque} \ill{} \ill{}
\note{les trois lignes qui précèdent l'encadré sont d'une écriture très rapide ; seuls les symboles se lisent sûrement.}
\[
  \begin{aligned}
    g_i\gamma_ix &= (g_j,g_k)x\\
    g_i(\gamma_i,\gamma_j,\gamma_k)x &= (\gamma_j,\gamma_k)\,g_i\gamma_ix = (\gamma_jg_j,\gamma_kg_k)x
  \end{aligned}
\]
\note{ces deux lignes sont encadrées dans la marge gauche.}
orbite, qui \ill{} \uncertain{est un} torseur sous l'un quelconque des \uncertain{trois} groupes $\prod_{j\in I\setminus\{i\}}G_j$. On a ainsi \struck{\uncertain{contenu}} \add{une} \uncertain{partie}
\[
  \Gamma\subset\prod_{i\in I}E_i=P .
\]
Celle-ci satisfait à \struck{\ill{}} \add{la} condition de \textbf{I}.
\[
  \begin{array}{lll}
    a=(a_i,a_j,a_k) & x_i=a_ig_k^{-1} & \text{détermine } g_k \text{ en fonction de } a \text{ et } x_i\\
    x=(x_i,x_j,x_k) & x_j=g_ka_jg_i^{-1} & \text{détermine } g_i \text{ en fonction de } a \text{ et } x_i,x_j\\
    \phantom{x=}\ ? & x_k=g_ia_k & \text{détermine } x_k
  \end{array}
\]

\page{5}
\note{la page commence au milieu d'une phrase, dont le début se lit au bas de la page 6 (« Donc on a trois applications $u_1,u_2,u_3 : G\to G'$ »).}
satisfaisant
\[
  u_i(1)=1 \qquad i=1,2,3
\]
et
\[
  (*)\qquad xyz=1\ \Longrightarrow\ u_1(x)\,u_2(y)\,u_3(z)=1
\]
Prouvons que les $u_i$ sont \struck{des homom.\ de} \add{égaux à un seul} et unique homom.\ de groupes.

Faisons $x=1$, on trouve que $u_3(y^{-1})=u_2(y)^{-1}$. Donc $u_3(y)=u_2(y^{-1})^{-1}$, pareil \ill{} : $u_1(y)=u_2(y^{-1})^{-1}$, donc $u_1(y)=u_3(y)$, et de même $u_1(y)=u_2(y)$. Donc $u_1=u_2=u_3$, et $u(y)=u(y^{-1})^{-1}$, donc \struck{$u$} $u(y^{-1})=u(y)^{-1}$.

On tire maintenant de $(*)$
\[
  u\bigl((xy)^{-1}\bigr)^{-1}=u(x)\,u(y),
\]
le membre de gauche valant $u(xy)$, donc $u(xy)=u(x)\,u(y)$ cqfd.

\emph{Question} : le foncteur $\varphi$ est-il une équivalence ?

\page{6}
Catégorie $\mathcal{C}'$ ($I$ variable) \uncertain{pour} iso sur $I$)

On a un foncteur
\[
  \mathcal{C}'\xrightarrow{\ \varphi\ }\mathcal{C}
\]

\emph{Prop.} Ce foncteur est pl.\ fidèle \add{(\uncertain{cet} identifie} $\mathcal{C}'$ à une sous-catégorie pleine de $\mathcal{C}$)
\note{la parenthèse est réécrite au-dessus d'un premier début biffé, illisible.}

Comme ce sont des catégories fibrées sur le groupoïde des ens.\ $I$ de card~$3$, on est ramené à prouver que \uncertain{le foncteur induit} sur les fibres pour $I=\{1,2,3\}$ \struck{\ill{}} \add{\uncertain{équivalentes}} est pl.\ fid. \struck{\ill{}} La \emph{fidélité} est claire $(E_i,G_i,\ldots,\Gamma)$.

Soit un objet de $\mathcal{C}'$, et $x$ \add{$=(x_1,x_2,x_3)$} un pt de $\Gamma$. Donc les $G_i$ sont isomorphes entre eux, et isomorphes aux $E_i$, et $\prod E_i$ s'identifie à $G^3$, et $\Gamma$ à l'ens.\ des $(x,y,z)$ tels que
\[
  xyz=1 .
\]
Soit un autre objet $(E'_i,G'_i,\ldots,\Gamma')$ et \ill{} \add{$\varphi$ \uncertain{dans}} \uncertain{un morphisme} \ill{} \ill{} $\varphi$ \ill{}, \uncertain{appliquant} $\Gamma$ dans $\Gamma'$, soit $(x'_1,x'_2,x'_3)$ l'image de $x=(x_1,x_2,x_3)=(1,1,1)$. Elle permet d'identifier le deuxième objet à $G'^3$ ($G'$ un groupe), $\Gamma'$ à l'ens.\ des $(x',y',z')$ tels que
\[
  x'y'z'=1 .
\]
Donc on a trois applications
\[
  u_1,u_2,u_3 : G\to G'
\]
\note{la phrase se poursuit en tête de la page 5 (« satisfaisant $u_i(1)=1$ … »), qui est donc la suite de celle-ci : les feuillets sont classés dans le désordre.}

\page{8}
\uncertain{$x\ y$} \qquad $a_1\ a_2\ a_3$
\[
  E_1\xrightarrow{\ \alpha_1\ }\mathrm{Isom}(E_2,E_3)\qquad
  E_2\xrightarrow{\ \alpha_2\ }\mathrm{Isom}(E_3,E_1)\qquad
  E_3\xrightarrow{\ \alpha_3\ }\mathrm{Isom}(E_1,E_2)
\]

\begin{tikzcd}[column sep=large, row sep=large]
E_1 \arrow[r, "\alpha_3(a_3)"] & E_2 \arrow[dl, "\alpha_1(a_1)"] \\
E_3 \arrow[u, "\alpha_2(a_2)"] &
\end{tikzcd}

\note{à l'intérieur du triangle, trois flèches en pointillé de sens inverse, marquées $u_3$ ($E_2\to E_1$), $u_1$, $u_2$.}

\emph{NB}
\[
  \alpha_3(a_3) : a_1\mapsto a_2,\qquad \alpha_1(a_1) : a_2\mapsto a_3,\qquad \alpha_2(a_2) : a_3\mapsto a_1
\]
Posons
\[
  \begin{aligned}
    \alpha_2(a_2)\,\alpha_1(a_1) &= u_3 : E_2\to E_1\\
    \alpha_3(a_3)\,\alpha_2(a_2) &= u_1 : E_3\to E_2\\
    \alpha_1(a_1)\,\alpha_3(a_3) &= u_2 : E_1\to E_3
  \end{aligned}
\]
\emph{Admettons}
\[
  \boxed{u_1u_2u_3=\mathrm{id}}
  \qquad \text{NB}\quad u_3(a_2)=a_1,\quad u_1(a_3)=a_2,\quad u_2(a_1)=a_3
\]
D'où syst.\ transitif d'isom.\ entre les $E_i$,
\note{la phrase se poursuit hors de la page.}

\page{10}
\struck{\ill{}} permettant de les identifier : \uncertain{à} un ens.\ $G$. On a donc
\note{suite, selon toute apparence, de la dernière ligne de la page 8 (« D'où syst.\ transitif d'isom.\ entre les $E_i$, »).}
\begin{itemize}
  \item ens.\ pointé $(G,e)$
  \item $\Gamma\subset G\times G\times G$
\end{itemize}
a) $\forall x,y\in G$,
\[
  \begin{aligned}
    &\exists!\,z\in G \text{ tel que } (x,y,z)\in\Gamma\\
    &\exists!\,z\in G \text{ tel que } (z,x,y)\in\Gamma\\
    &\exists!\,z\in G \text{ tel que } (x,z,y)\in\Gamma
  \end{aligned}
\]
\struck{b) Soit \ill{} $\sigma_1,\sigma_3$ : $G\xrightarrow{\sim}G$ défini par $(x,\sigma x,e)\in\Gamma$, $(e,x,\sigma_1x)\in\Gamma$, $(\sigma_2x,e,x)\in\Gamma$}
\note{ce premier b) est barré de hachures obliques et encadré.}

b) $(e,e,e)\in\Gamma$

c) Soit $\sigma_1,\sigma_2,\sigma_3 : G\xrightarrow{\sim}G$ défini par
\[
  \begin{aligned}
    &(e,x,\sigma_1x)\in\Gamma\\
    &(\sigma_2x,e,x)\in\Gamma\\
    &(x,\sigma_3x,e)\in\Gamma
  \end{aligned}
\]
On a \struck{$(\sigma_3\sigma_2\sigma_1)^2=\mathrm{id}$}
\[
  \sigma_2\sigma_1=\mathrm{id},\quad \sigma_3\sigma_2=\mathrm{id},\quad \sigma_1\sigma_3=\mathrm{id}
\]
i.e.\ $\sigma_1=\sigma_2=\sigma_3$ $(=\sigma)$, $\sigma^2=\mathrm{id}$.

Si $x,y\in G$ soit $x.y\in G$ défini par $(x,y,\sigma(x.y))\in\Gamma$. On a donc \struck{$z=xy\Longleftrightarrow(x,y$}
\[
  (x,y,z)\in\Gamma\ \Longleftrightarrow\ z=\sigma(x.y)
\]

\page{12}
$(G,e)$ ens.\ pointé, avec loi de composition $(x,y)\mapsto x.y$, $\sigma$ automorphisme de $G$ \struck{$e$} avec $\sigma^2=\mathrm{id}$
\[
  \Gamma=\{(x,y,z)\in G^3 \mid z=\sigma(x.y)\}=\{(x,y,z)\in G^3\mid x.y=\sigma^{-1}z\}
\]
\emph{Conditions} a) $\forall x,y\in G$
\[
  \begin{aligned}
    &\exists!\,z \text{ tel que } zx=\sigma^{-1}y\\
    &\exists!\,z \text{ tel que } xz=\sigma^{-1}y
  \end{aligned}
\]
i.e.\ $z\mapsto zx$ et $z\mapsto xz$ sont bijectifs

b) $\sigma^{-1}e=e.e$

c) On a
\[
  \begin{aligned}
    &\forall x\quad \sigma^{-1}\sigma x=ex &&\text{i.e. } ex=x\\
    &\forall x\quad \sigma x.e=\sigma^{-1}x &&\text{i.e. } \forall x \text{ on a } x.e=\sigma^{-2}x\ (=x)\\
    &\forall x\quad x.\sigma x=\sigma^{-1}e\ (=e)
  \end{aligned}
\]
\note{en marge de b) et de la première ligne de c), une accolade conclut « $\Rightarrow\sigma e=e$ ».}

\struck{y faisant $x=1$, on trouve $e.\sigma e=\sigma^{-1}e$ $(=\sigma e)$ donc $e.\sigma$}

Donc : $e$ est unité bilatère, \struck{et $\sigma x$ est l'inverse à droite de $x$} et les translations à g.\ et à dr.\ sont bijectives. \uncertain{De plus}, $\forall x\in G$, $\sigma x$ est \struck{\uncertain{inverse}} \add{\uncertain{inverse}} à droite \add{de $x$}, donc \uncertain{(comme $\sigma^2=\mathrm{id}$)}
\[
  x.\sigma x=e\qquad (\sigma x)(\sigma\sigma x)=e
\]
\struck{$(\sigma^{-1}x)x=e$} \quad c'est un inverse bilatère \struck{$\sigma x.x$}
\note{l'alinéa « Donc : $e$ est unité bilatère … » est marqué d'un double trait vertical en marge gauche.}

\[
  (xy)z=x(yz)\ \ ??\qquad (x,y,\sigma(xy))\in\Gamma
\]
\struck{$(y\sigma x$} \quad \struck{\ill{}} \quad \struck{$xy,z,\sigma((xy)z)$}
\[
  (x,y,z)\in\Gamma\ \Longrightarrow\ z=\sigma(x.y)\ \Longleftrightarrow\ (x.y).z=e\ \Longleftrightarrow\ z.(x.y)=e
\]
\[
  \begin{aligned}
    &(xy,z,\sigma((xy)z))\in\Gamma\\
    &(x,yz,\sigma(x(yz)))\in\Gamma
  \end{aligned}
  \qquad
  \begin{aligned}
    &xy=e\\
    &\Updownarrow\\
    &(xy)z=z\\
    &\Updownarrow\\
    &z(xy)=z
  \end{aligned}
\]
\note{le passage de « $z=\sigma(x.y)$ » à « $(x.y).z=e$ » est écrit au-dessus d'un essai biffé, \struck{$(x.y)=\sigma x$}.}

\page{13}
\marginal{(A)}
\[
  E_1=E_2=E_3=G\qquad \Gamma=\{(x,y,z)\mid xyz=1\}
\]
$(a,b,c)\in\Gamma$ \quad i.e.\ $abc=1$

\begin{tikzcd}[column sep=large, row sep=large]
& E_1 \arrow[dl, "f_3"'] & \\
E_2 \arrow[rr, "f_1"] & & E_3 \arrow[ul, "f_2"'] \arrow[ll, bend left=30, "g_1"]
\end{tikzcd}

\note{sur le schéma, $g_3$ et $g_2$ sont des flèches courbes parallèles à $f_3$ et $f_2$, de sens contraire.}
\[
  \begin{aligned}
    &f_1(x) \text{ défini par } (a,x,f_1(x))\in\Gamma\\
    &f_2(x) \text{ défini par } (f_2(x),b,x)\in\Gamma\\
    &f_3(x) \text{ défini par } (x,f_3(x),c)\in\Gamma
  \end{aligned}
\]
\[
  f_1(x)=x^{-1}a^{-1},\qquad f_2(x)=x^{-1}b^{-1},\qquad f_3(x)=x^{-1}c^{-1}
\]
\[
  x^{-1}a^{-1}=y,\qquad x=a^{-1}y^{-1}
\]
\[
  \boxed{\sigma(xy)=\sigma y\,\sigma x}
\]
\[
  \text{\struck{$z=xy$}}\qquad
  \left\{
  \begin{aligned}
    &(xy)z=1\\ &\Updownarrow\\ &(zx)y=1\\ &\Updownarrow\\ &(x^{-1}z^{-1})y^{-1}=1
  \end{aligned}
  \right.
  \qquad
  \begin{aligned}
    &z=xy \Longleftrightarrow x=zy^{-1}\\
    &y^{-1}=zx,\quad y=x^{-1}z^{-1}
  \end{aligned}
\]
\note{ce coin droit est très surchargé (plusieurs essais biffés) et séparé du reste par un trait ; lecture d'ensemble douteuse.}
\[
  \begin{aligned}
    g_1(x)&=f_3f_2(x)=(bx)c^{-1}\\
    g_2(x)&=f_1f_3(x)=(cx)a^{-1}\\
    g_3(x)&=f_2f_1(x)=(ax)b^{-1}
  \end{aligned}
\]
\[
  g_1g_2g_3(x)=bcaxb^{-1}a^{-1}c^{-1}=(bca)\,x\,(cab)^{-1}=x\qquad\text{i.e. } g_1g_2g_3=\mathrm{id}
\]
\note{sous $(bca)$ et $(cab)$, deux accolades marquées « $1$ ».}
\[
  g_1g_2g_3=(f_3f_2f_1)^2
\]
\[
  f_3f_2f_1(x)=(axb^{-1})^{-1}c^{-1}=bx^{-1}a^{-1}c^{-1}=bx^{-1}(ca)^{-1}=bx^{-1}b
\]
\[
  \sigma_2x=bx^{-1}b\qquad\text{de même}\qquad \sigma_3x=cx^{-1}c,\quad \sigma_1x=ax^{-1}a
\]
\struck{$\sigma_2^2x=b(b^{-1}xb^{-1})b=x$}

\[
  E_1\qquad E_2\ \underset{f_1^{-1}}{\overset{f_1}{\rightleftarrows}}\ E_3
  \qquad \text{\struck{$(x,y,z)$}}
\]
\note{à droite, un petit schéma : $\sigma_1$ de $E_1$ dans $E_1$, et $f_1$, $f_1^{-1}$ échangeant $E_2$ et $E_3$ en croix.}
\[
  \begin{aligned}
    &(x,y,z)\\
    &(\sigma_1x,\ f_1^{-1}z,\ f_1y)=(ax^{-1}a,\ a^{-1}z^{-1},\ y^{-1}a^{-1})
  \end{aligned}
  \qquad
  \begin{aligned}
    &xyz=1\Longrightarrow\\
    &ax^{-1}aa^{-1}z^{-1}y^{-1}a^{-1}=1
  \end{aligned}
\]

\page{14}
\marginal{(B)}
\note{la page reprend le calcul de la page 13 sans supposer l'associativité : les parenthèses sont écrites partout.}
\struck{$z=xy$} \qquad $x^{-1}(uz^{-1})=y$
\[
  \begin{aligned}
    u=(xy)z &\Longleftrightarrow (uz^{-1}=xy)\Longleftrightarrow (y^{-1}x^{-1})u=z\\
    u=x(yz) &\Longleftrightarrow (x^{-1}u=yz)\Longleftrightarrow u(\cdots
  \end{aligned}
\]
\note{la seconde ligne s'interrompt sur « $u($ » ; suivent, sous un $x^{-1}$ isolé, quelques essais biffés illisibles.}
\[
  \left\{
  \begin{aligned}
    f_1(x)&=x^{-1}a^{-1}\\ f_2(x)&=x^{-1}b^{-1}\\ f_3(x)&=x^{-1}c^{-1}
  \end{aligned}
  \right.
  \qquad
  \begin{aligned}
    g_1&=f_3f_2 = x\mapsto (bx)c^{-1}\\
    g_2&=f_1f_3 = x\mapsto (cx)a^{-1}\\
    g_3&=f_2f_1 = x\mapsto (ax)b^{-1}
  \end{aligned}
\]
\[
  g_1g_2g_3=(f_3f_2f_1)^2\qquad x\mapsto \Bigl(b\Bigl(\bigl(c((ax)b^{-1})\bigr)a^{-1}\Bigr)\Bigr)c^{-1}
\]
\[
  \sigma_2x=f_3f_2f_1(x)=\bigl((ax)b^{-1}\bigr)^{-1}c^{-1}=\bigl(b(x^{-1}a^{-1})\bigr)c^{-1}
\]
\[
  \sigma_2^2x=\Bigl(b\bigl(\bigl[c((ax)b^{-1})\bigr]a^{-1}\bigr)\Bigr)c^{-1}
\]
\[
  \begin{array}{l}
    x\\ ax\\ axb^{-1}\\ caxb^{-1}\\ \text{\struck{$b$}}\,caxb^{-1}a^{-1}\\ bcaxb^{-1}\\ bcaxb^{-1}a^{-1}c^{-1}
  \end{array}
  \qquad
  \left\{
  \begin{aligned}
    &\sigma(xy)=\sigma x\,\sigma y\\
    &\tau^{d}_{x}{}^{-1}=\tau^{d}_{x^{-1}}\\
    &\tau^{g}_{x}{}^{-1}=\tau^{g}_{x^{-1}}
  \end{aligned}
  \right.
\]
\note{$\tau^{d}$, $\tau^{g}$ : translations à droite et à gauche, lecture probable des exposants $d$, $g$. La colonne de gauche est la suite des étapes de $g_1g_2g_3x$ ; une ligne y est en partie biffée.}
\[
  u=(xy)z\qquad (bx)c^{-1}=y\qquad x=b^{-1}(yc)
\]
\[
  x=\qquad g_1g_2g_3x=x\qquad \text{\struck{$g_2g_3x=$}}
\]
\[
  \bigl(c((ax)b^{-1})\bigr)\text{\struck{$a^{-1}$}}=\bigl(b^{-1}(xc)\bigr)a
\]
\[
  \boxed{c\bigl((ax)b^{-1}\bigr)=\bigl(b^{-1}(xc)\bigr)a}
\]
\note{en marge gauche : « $abc$, $cab=1$ » et deux lignes biffées ; au pied de la page, deux essais biffés :}
\struck{$c(xc)=$ \ill{} $(ax)a=$ \ill{}} \struck{$c(ab^{-1})=(b^{-1}c)a$}

\page{15}
\note{page de calcul très surchargée, en deux encres (bleue et brune) ; seules les formules principales sont transcrites, les nombreux essais biffés et les accolades de soulignement sont résumés.}
\[
  (b^{-1}a^{-1})\bigl((ax)b^{-1}\bigr)=\bigl(b^{-1}(x(b^{-1}a^{-1}))\bigr)a
\]
\note{formule entourée, à l'encre brune.}
\[
  \left\{
  \begin{aligned}
    &(ab)c=1\\
    &\text{\struck{$x\ y\ z$}}\\
    &(xy)z=1
  \end{aligned}
  \right.
  \ \Longrightarrow\
  \Bigl[\bigl((z^{-1}b^{-1})((b(y^{-1}a^{-1}))c^{-1})\bigr)\Bigr](b^{-1}x^{-1})\,\text{\struck{$b$}}=1
\]
\note{ce bloc est encadré.}

\struck{$\bigl(b(x^{-1}a^{-1})\bigr)c^{-1}=b(x^{-1}b)$}
$z=c^{\cdot}$ \ill{}
\note{l'exposant de $c$ est illisible.}
\[
  b\mapsto z,\quad a\mapsto b,\quad x\mapsto y^{-1}a^{-1}
\]
\[
  \Bigl[\bigl(z^{-1}((y^{-1}a^{-1})(z^{-1}b^{-1}))\bigr)b\Bigr](b^{-1}x^{-1})
\]
\[
  \text{\struck{$c=b^{-1}a^{-1}$}}\quad \text{\struck{$z=ab$}}\qquad (ab)(xy)=1
\]
\[
  ab=xy\qquad \text{\struck{$(ab)(xy)=1$}}\ \Longrightarrow\
  \Bigl[(\cdots)\bigl((y^{-1}a^{-1})(\cdots\,b^{-1})\bigr)b\Bigr](b^{-1}x^{-1})
\]
\note{les deux facteurs notés ici $(\cdots)$ sont surchargés et illisibles.}
\note{sous la dernière expression, des accolades regroupent des facteurs en « $a$ », « $y^{-1}$ », « $x$ » ; plus bas : « $ax$ ».}
\[
  (xb)(b^{-1}x^{-1})=1
\]
\[
  \begin{aligned}
    &(ab)c=1\ *\\ &(xy)z=1\ *
  \end{aligned}
  \ \overset{?}{\Longrightarrow}\ xb=(z^{-1}b^{-1})\bigl((b(y^{-1}a^{-1}))c^{-1}\bigr)
\]
\[
  \Updownarrow
\]
\[
  (bz)(xb)=\bigl(b(y^{-1}a^{-1})\bigr)c^{-1}
\]
\struck{$(ub^{-1})(bc)=uc\ ?$}
\note{formule biffée et entourée, à l'encre brune, en marge gauche.}
\[
  (bz)(xb)=\bigl(b((zx)(bc))\bigr)c^{-1}
\]
\[
  \bigl((bz)(xb)\bigr)c=b(y^{-1}a^{-1})\qquad (bu)c=b\bigl((ub^{-1})(bc)\bigr)
\]
\note{sous $\bigl((bz)(xb)\bigr)c$, des accolades : $bz\mapsto b$, $xb\mapsto y^{-1}$, $c=a^{-1}$.}
\[
  \boxed{\bigl((bz)(xb)\bigr)c=b\bigl((zx)(bc)\bigr)}
  \qquad
  \left\{
  \begin{aligned}
    &c=a^{-1}\\ &xb=y^{-1}\ *\\ &b=bz\ *
  \end{aligned}
  \right.
\]
\note{la formule encadrée est à l'encre brune.}
\[
  z=1,\ c=\quad \bigl(b(xb)\bigr)c=b\bigl(x(bc)\bigr)\qquad z=1,\ y=x^{-1},\quad b=1,\ c=a^{-1}
\]
\note{dans les deux dernières égalités, $x$ et $a$ sont entourés ; lecture d'ensemble de cette ligne douteuse.}

\page{16}
\note{page entière barrée de deux grandes diagonales en croix. Elle porte deux couches d'écriture : à l'encre bleue, les essais transcrits ci-dessous ; tête-bêche, à l'encre brune, quelques formules (« $z\,(u(u^{-1}z^{-1}))$ », « $(xy)(y^{-1}x^{-1})$ », « $(aa)a$ ») qui reprennent les calculs d'associativité des pages précédentes.}
\[
  \begin{aligned}
    &\tau_xz=xz\\ &\tau_x\tau_yz=x(yz)\\ &xx'=x'x=e\\ &x(x'y)=y\\ &G\times G\to G
  \end{aligned}
\]
\emph{Petits ensembles.} $e\ a$ :\quad $ee=e$, $ea=ae=a$, $aa\neq ae$ donc $aa=e$.

$e\ a\ b$ :
\[
  aa\quad bb\quad ab=ba \qquad ab\overset{?}{=}ba,\quad ab=a(aa),\ ba=(aa)a
\]
\note{sous « $aa\ bb\ ab=ba$ », une ligne d'essais biffés en zigzag, d'où émergent « $b$ » et « $e$ ».}
\[
  \{ae,aa,ab\}=\{e,a,b\}\qquad aa=b,\ ab=e\ \text{ou}\ aa=e,\ \text{\struck{$ab=b$}}\ \text{mais}\ ab\neq\ill{}
\]
\note{le dernier terme, surchargé, est illisible.}
\[
  aa=b,\ ab=e\ \Longrightarrow\ \boxed{ba=e}\qquad bb=a
\]
\[
  e,a,b,c \qquad \text{\struck{$(xy)z=$}}\ x(yz)\qquad uv=u'v'
\]
\[
  x\,y\,z\qquad \bar x\bigl((xy)z\bigr)=(yz)\qquad \bar x\bigl((xy)z\bigr)z
\]
$e$, $a,b,c$ : deux tables de multiplication commencées,
\[
  \begin{array}{c|ccc}
    & a & b & c\\ \hline
    a & e & c & b\\ b & & & \\ c & & &
  \end{array}
  \qquad\qquad
  \begin{array}{c|ccc}
    & a & b & c\\ \hline
    a & b & c & e\\ b & & & \\ c & & &
  \end{array}
\]
\note{dans la première table, l'entrée $(a,b)$ est récrite sur une autre lettre.}
\[
  aa=e\quad ab=c\quad ac=b
\]
\[
  (xy)z=1\ \Longrightarrow\ z(xy)=1\qquad \Updownarrow\qquad y(zx)=1,\quad (zx)y=1\qquad \text{\struck{$(xy)z$}}
\]
\[
  \begin{aligned}
    &(ab)c=1\\ &(ax)y=1\\ &(yb)x=1\\ &(xy)c=1
  \end{aligned}
  \qquad
  \begin{aligned}
    &y=\sigma(ax)\\ &yb=\sigma x\\ &xy=\sigma c
  \end{aligned}
  \qquad
  \begin{aligned}
    &f_1(x)=\sigma(ax)\\ &f_2(x)=(\sigma x)b^{-1}
  \end{aligned}
\]

\section{Un graphe construit sur un groupe}

\page{18}
\note{nouvelle construction, sur un feuillet blanc ; elle n'est pas reliée explicitement à ce qui précède.}
\begin{itemize}
  \item $V$ groupe
  \item $I\subset V$ partie de $V$
  \item $V\xrightarrow{\ \varphi_\lambda\ }V_\lambda$ ($\lambda\in\Lambda$) homom.\ de groupes, $I_\lambda\subset V_\lambda$ partie de $V_\lambda$
  \item \struck{$X$} \add{$Y$} un $V$-ens. \quad $X_\lambda=Y\wedge_V(V_\lambda,\varphi_\lambda)$, d'où $t_\lambda : Y\to X_\lambda$
\end{itemize}
\note{au-dessus du « $Y$ » de $X_\lambda=Y\wedge_V(\ldots)$, un « $X$ » rattaché par un trait : la lettre $X$ d'abord écrite a été remplacée par $Y$.}
\[
  S=\text{\struck{$X$}}\ \{e\}\sqcup Y\sqcup\coprod_{\lambda\in\Lambda}\underbrace{Y\times_V(V_\lambda,\varphi_\lambda)}_{X_\lambda}
\]
Structure de graphe sur $S$ :
\begin{itemize}
  \item $e$ est lié à tous les él.\ de \struck{\ill{}} $X$\ill{} ;
  \item \struck{les} $y\in Y$ est lié à $iy$ ($i\in I$)
  \item $x\in X_\lambda$ est lié à $ix$ ($i\in V_\lambda$)
  \item $y\in Y$ est lié à son image $t_\lambda(y)\in X_\lambda$
\end{itemize}
et cela épuise les arêtes.
\note{au premier item, le complément après « de » est surchargé, et l'indice de $X$ illisible. Au troisième, la page porte « $i\in V_\lambda$ » ; on attendrait plutôt $i\in I_\lambda$.}

À quelles conditions est-ce là un graphe \uncertain{triadique} ? Je dis qu'il n'y a que les cas suivants

a) $V\simeq\mathbf{F}_2^{\,r+1}$ \struck{\ill{}}, \uncertain{où} card \ill{} $=2^r$ \struck{\ill{}}, $r\in\{0,1,2\}$

b) $I$ est une base de $V$

c)
\note{la liste s'arrête à « c) », laissé vide.}

\page{20}
\drawing{En haut à gauche, un hexagone dont les six sommets portent des triplets de lettres $\xi$ à indices (« $\xi_1\xi_2\xi_3$ », « $\xi_4\xi_5\xi_6$ », … ; certains indices surchargés, d'autres marqués d'un accent ou d'un trait), reliés par les côtés en trait plein et par des diagonales en pointillé qui se croisent au centre ; une flèche courbe mène vers la droite.}
\note{les indices des triplets aux sommets de l'hexagone ne se lisent pas tous ; ils ne sont pas transcrits.}

$E^{\varepsilon}$
\drawing{À droite de $E^{\varepsilon}$, une rangée de points encadrée dont les deux derniers sont reliés à un point inférieur en forme de V, et sous elle une rangée de points pointillée terminée par un segment vertical.}

(\uncertain{hexagone} \uncertain{distingué} \uncertain{ou triplets typiques})
\note{sous la parenthèse, un « $\xi$ » isolé suivi de deux chiffres peu lisibles (« $1\ 4$ » ?).}

\struck{$2$ ens : trois él.\ $E$, $E'$}

$1$ ens : trois él.\ $I$
\[
  E\times I\ \sqcup\ E\times I'
\]
ens : $E$ él.\ fibrés \add{de degré $3$} sur ens $\varepsilon$ : $2$ él.
\[
  E=\coprod_{\alpha\in\varepsilon}E_\alpha
\]
ens : trois él.\ $I$
\note{ces trois lignes sont réunies par une accolade ; le « $E$ » de la première est récrit sur un autre signe.}
\[
  I\times\varepsilon\quad \text{hexagone}
\]
\[
  \coprod_{\alpha\in\varepsilon}E_\alpha\times I = E\times I \quad\text{fibré sur } I\times\varepsilon
\]

$\xi_1\xi_2\xi_3$ \quad il y a $6$ él.\ liés à aucun \struck{\uncertain{autre}} \ill{}

il y a \uncertain{trois} él.\ liés à deux \ill{} $\alpha\in E$

il y a $6$ él.\ liés à exactement $2$ des $\xi$

et $9$ él.\ liés à exactement un
\[
  X=\prod_{\alpha\in\varepsilon}E_\alpha
\]
\drawing{Le bas de la page est couvert d'esquisses sans légende : colonnes et grilles $3\times 3$ de points (l'une entourée d'un cercle), un faisceau de segments croisés entre deux rangées de points, de petits polyèdres, deux triangles, et les inscriptions « $x\,y\,z\,x'y'z'$ » (avec une accolade), « $9\ 6\ 6\ 6$ », « $9\ 9\ 9$ » et « $40$ ». Dans la marge droite, quelques traits illisibles.}
\note{la page s'arrête sur ces esquisses ; l'argument se poursuit vraisemblablement au-delà du lot.}

\end{document}
